\documentclass[10pt,a4paper]{amsart}
\usepackage[margin=19mm]{geometry}
\usepackage{amsmath,amssymb,mathtools}
\usepackage[scr=rsfs,cal=euler]{mathalfa}
\usepackage{xfrac}
\usepackage[unicode,hidelinks,bookmarks=false]{hyperref}
\newcommand{\ie}{i.e.\ }
\newcommand{\cf}{cf.\ }
\newcommand{\resp}{respectively\ }
\newcommand{\Subsection}{\subsection}
\providecommand{\nobreakdash}{\nobreak\hbox{-}}
\setlength{\emergencystretch}{2em}
\multlinegap0pt
\usepackage{amssymb}
\usepackage{tikz}
\usepackage{float}
\usepackage{xcolor}
\usepackage[normalem]{ulem}
\newtheorem{claim}{Claim}
\newcommand{\conv}{\mathsf{conv}} 
\renewcommand{\subset}{\subseteq}
\title{Infima and cardinal characteristics of critical ideals for countable compact spaces}
\author{Malgorzata Kowalczuk}
\date{Source: arXiv:2603.00674v1 (CC BY 4.0)}
\begin{document}
\maketitle

\noindent\emph{Source and licence.} Infima and cardinal characteristics of critical ideals for countable compact spaces. Version arXiv:2603.00674v1 (CC BY 4.0), distributed by arXiv under the Creative Commons Attribution 4.0 International licence, http://creativecommons.org/licenses/by/4.0/, as recorded in the arXiv licence metadata for this exact version and shown on the abstract page https://arxiv.org/abs/2603.00674v1. Full licence notice: \texttt{licenses/arxiv-2603.00674-CC-BY-4.0.txt}.\par\smallskip

\noindent\textit{Editorial scope.} Counted unit: the article's claim claim (source0.tex lines 482--484). The article's complete original proof of it is delivered with it (source0.tex lines 486--514, one \texttt{proof} environment of 29 lines). No mathematics is written, repaired or re-derived: the statement and the proof are the article's own lines, in the article's own notation, and no retained line is substituted or re-flowed.  No further source-local context is needed: every cross-reference of every retained line resolves inside the two delivered spans.  Nothing was omitted from any retained span, so the package carries no omission record.  No retained line contains a citation, so the wrapper carries no bibliography and no bibliography entry.  No retained line contains a figure environment, an image-inclusion command or drawing code, so no image is delivered and the package declares no asset dependency.  Editorial formatting only, in addition: an AMS \texttt{amsart} preamble that restates the article's own declarations and macros used by the retained lines, namely the environments \texttt{claim} on separate counters, together with the standard packages those lines need.  The retained environments print in this document as Claim 1 in order of appearance, whereas the article prints the same items with the numbering of its own full document.  The complete native source of the article as distributed in the arXiv e-print for this exact version is bundled unchanged as \texttt{sources/195497/source0.tex}.
\begin{claim}    
There is some $k_{i+1}>\max(n_i,k_i)$ such that the set $$(f[[\omega^{\alpha_{k_{i+1}}},\omega^{\alpha_{k_{i+1}}}+\omega^{\alpha_{n_i}}]])^d\setminus (\omega^{\alpha_{n_i}}+1)$$ is infinite.
\end{claim}

\begin{proof}[Proof of Claim]
Suppose for the sake of contradiction that the set $(f[[\omega^{\alpha_k},\omega^{\alpha_k}+\omega^{\alpha_{n_i}}]])^d\setminus (\omega^{\alpha_{n_i}}+1)$ is finite for every  $k>\max(n_i,k_i)$.
For every such $k$, define the set $$F_k=f[[\omega^{\alpha_k},\omega^{\alpha_k}+\omega^{\alpha_{n_i}}]]\setminus (\omega^{\alpha_{n_i}}+1)$$ and note that since $(F_k)^d$ is finite, $F_k\in\conv_{<\alpha}$, and therefore $f^{-1}[F_k]\in\conv_\alpha$.  Let $g: \bigcup_{k>\max(n_i,k_i)} [\omega^{\alpha_k},\omega^{\alpha_k}+\omega^{\alpha_{n_i}}]\to\omega^{\alpha_{n_i}}+1$ be a function such that 

\begin{enumerate}
    \item $g$ coincides with $f$ on the set
$\bigcup_{k>\max(n_i,k_i)} [\omega^{\alpha_k},
 \omega^{\alpha_k}+\omega^{\alpha_{n_i}}]\setminus f^{-1}[F_k]$,


 \item $g\restriction ([\omega^{\alpha_k},\omega^{\alpha_k}+\omega^{\alpha_{n_i}}]\cap f^{-1}[F_k])$ is injective for each $k>\max(n_i,k_i)$,
 \item $g[ [\omega^{\alpha_k},\omega^{\alpha_k}+\omega^{\alpha_{n_i}}]\cap f^{-1}[F_k]]\subset[\omega\cdot k, \omega\cdot (k+1)).$ 

 \end{enumerate}

We will show that $g $ witnesses $\conv_{\alpha_{n_i}}\leq_K \conv(\bigcup_{k>\max(n_i,k_i)} [\omega^{\alpha_k},\omega^{\alpha_k}+\omega^{\alpha_{n_i}}])$ which will give us a contradiction since 
$$\conv\left(\bigcup_{k>\max(n_i,k_i)} [\omega^{\alpha_k},\omega^{\alpha_k}+\omega^{\alpha_{n_i}}]\right) \approx \conv_{\alpha_{n_{i}}+1}\not\geq_K\conv_{\alpha_{n_i}}.$$

To see that the function $g$ witnesses the Kat\v{e}tov  reduction, take any $A\notin  \conv(\bigcup_{k>\max(n_i,k_i)} [\omega^{\alpha_k},\omega^{\alpha_k}+\omega^{\alpha_{n_i}}])$. Since $A^d$ is infinite we have two cases.

\emph{Case 1.} There is  $k>\max(n_i,k_i)$ such that $A\cap [\omega^{\alpha_k},\omega^{\alpha_k}+\omega^{\alpha_{n_i}}]\notin \conv([\omega^{\alpha_k},\omega^{\alpha_k}+\omega^{\alpha_{n_i}}])$. Then $B=A\cap [\omega^{\alpha_k},\omega^{\alpha_k}+\omega^{\alpha_{n_i}}]\setminus f^{-1} [F_k]\notin \conv([\omega^{\alpha_k},\omega^{\alpha_k}+\omega^{\alpha_{n_i}}]) $, so $g[B]=f[B]\notin \conv_{<\alpha}$. However, $g[B]\subset [0,\omega^{\alpha_{n_i}}]$ and therefore $g[B]\notin\conv_{\alpha_{n_i}}.$

\emph{Case 2.} The set $\{l>\max(n_i,k_i) : |A\cap [\omega^{\alpha_l}, \omega^{\alpha_l}+\omega^{\alpha_{n_i}}]|=\omega\}$ is infinite.
Let $\{l_j :j\in\omega\}$ be an
enumeration of the set above, and let $B_l=A\cap [\omega^{\alpha_l}, \omega^{\alpha_l}+\omega^{\alpha_{n_i}}]$ for every $l$.
Then
$$g\left[\bigcup_{j\in\omega} B_{l_j}\right]=f\left[\bigcup_{j\in\omega} B_{l_j}\setminus f^{-1}[F_{l_j}]\right]\cup g\left[\bigcup_{j\in\omega} B_{l_j}\cap f^{-1}[F_{l_j} ]\right] .$$
Since $\bigcup_{j\in\omega} B_{l_j}\notin\conv_\alpha$, either $\bigcup_{j\in\omega} B_{l_j}\setminus f^{-1}[F_{l_j}]\notin\conv_\alpha$ or  $\bigcup_{j\in\omega} B_{l_j}\cap f^{-1}[F_{l_j} ]\not\in\conv_\alpha$. In the former case, it follows that $f[\bigcup_{j\in\omega} B_{l_j}\setminus f^{-1}[F_{l_j}]]\notin\conv_{<\alpha}$. But $f[\bigcup_{j\in\omega} B_{l_j}\setminus f^{-1}[F_{l_j}]]\subset [0,\omega^{\alpha_{n_i}}]$ and therefore $g[A]\supseteq f[\bigcup_{j\in\omega} B_{l_j}\setminus f^{-1}[F_{l_j}]]\notin\conv_{\alpha_{n_i}}$. In the latter case, let us note that since $f^{-1}[F_{l_j}]\in\conv_\alpha$ for each $j$, the set $\{j\in\omega : |B_{l_j}\cap f^{-1}[F_{l_j}]|=\omega\}$ is infinite. By the definition of $g$, we get that $ g[\bigcup_{j\in\omega} B_{l_j}\cap f^{-1}[F_{l_j} ]] \notin\conv_{\alpha_{n_i}}$.
\end{proof}

\end{document}
